\section{Wstęp i założenia}\label{sec:wstep}
Celem pracy jest wykazanie, że każdy dostatecznie silny, efektywnie aksjomatyzowalny i niesprzeczny system formalny arytmetyki jest \emph{niezupełny}: istnieje zdanie arytmetyczne, którego system ten ani nie dowodzi, ani nie obala. Ponadto pokażemy, że taki system nie może dowieść własnej niesprzeczności.

W pierwszym twierdzeniu pracujemy z teorią $T$ spełniającą warunki:
\begin{enumerate}[label=(\roman*)]
\item językiem jest język arytmetyki pierwszego rzędu, a rachunkiem logicznym --- ustalony klasyczny rachunek dowodowy;
\item aksjomaty pozalogiczne można \emph{efektywnie wyliczać}: ich zbiór jest rekurencyjnie przeliczalny;
\item $T$ zawiera arytmetykę Robinsona $Q$;
\item $T$ jest niesprzeczna. Dla obustronnej niezależności konkretnego zdania Gödla zakładamy mocniejszą $\omega$-niesprzeczność; dla zdania Rossera wystarcza zwykła niesprzeczność.
\end{enumerate}
Rozstrzygalność zbioru aksjomatów jest warunkiem wystarczającym, lecz nie koniecznym. W dowodzie drugiego twierdzenia i warunków wyprowadzalności przyjmujemy mocniejszą bazę $T\supseteq\mathrm{PA}$ i ściśle określony, standardowy predykat dowodliwości. Nie przenosimy na $Q$ argumentów wymagających indukcji.

Rozróżniamy metateorię, w której opisujemy skończone dowody, od zdań dowodzonych \emph{wewnątrz} $T$. Zwłaszcza rodzina twierdzeń $T\vdash\theta(\overline n)$ dla wszystkich standardowych $n$ nie oznacza automatycznie $T\vdash\forall x\,\theta(x)$.

